\begin{textsample}{POS Dim 3 – to – Score 21.00 – to/to\_inf\_31.txt}  \label{ex:f3_pos_110}
Campo de experiência : espaços, tempos, quantidades, relações e transformações As \textbf{crianças} são sujeitos ativos, curiosos e observadores, vivem inseridas em um mundo formado por diversos fenômenos naturais, culturais e sociais.
Buscam compreender os diversos espaços que convivem, os tempos, as quantidades, as relações e as transformações desse ambiente.
Desde \textbf{bebês}, as \textbf{crianças} são curiosas, pois buscam compreender o mundo que as cerca.
Elas possuem grande interesse em descobrir como as coisas são, suas transformações e comportamentos.
Buscam entender o seu próprio corpo, como ele se expressa, os espaços que \textbf{conseguem} ocupar, os tempos que utilizam para realizar \textbf{brincadeiras} e dançar músicas.
Vivências como essas, contribuem para o desenvolvimento da compreensão corporal.
Os fenômenos da natureza, quer sejam os naturais, físicos ou biológicos, também provocam interesse nas \textbf{crianças}.
O processo de mediação do professor é essencial durante o desenvolvimento da \textbf{criança}, pois, ao manifestar a \textbf{curiosidade} em compreender uma ação ou fenômeno, o educador precisa alimentar o \textbf{desejo}, e permitir a aproximação das \textbf{crianças} nessas experiências ricas e significativas.
Quando já \textbf{conseguem} falar, gostam de fazer perguntas sobre tudo que lhes desperta interesse, “ a pergunta mais constante que as \textbf{crianças} \textbf{pequenas} expressam com olhares interessados, é : “ Por quê? ” E não adianta dizer : “ Porque sim ” ou “ Porque não ”, nem desfilar conceitos e teorias diante das \textbf{crianças}.
Na potente caminhada de cada \textbf{criança} para produzir saberes, o desafio maior é entender e responder às falas \textbf{infantis}, perceber as relações que as \textbf{crianças} estabelecem entre fatos, descobrir as teorias que elaboram, o que requer incentivar que elas façam perguntas, sejam mais curiosas ( BRASIL. p. 94, 2018 ).
Nesse sentido, os professores precisam oferecer às \textbf{crianças} oportunidades para investigarem os diversos assuntos e objetos do seu cotidiano.
Assim, elas formulam questões, levantam hipóteses, \textbf{conseguem} respostas e aprendem de maneira significativa.
A construção social de conhecimentos pela \textbf{criança} \textbf{pequena} depende das situações criadas pelos \textbf{parceiros} mais experientes para mediar suas aprendizagens.
Em outras palavras, nas interações cotidianas os \textbf{parceiros} emprestam à \textbf{criança} sua forma de selecionar e relacionar elementos, seu modo de explicar algo, recorrendo a uma fala que inclui descrições ( “ isto é azedo ” ), ou hipóteses ( “ se o \textbf{carrinho} não anda é porque sua pilha acabou ” ), até que a própria \textbf{criança} crie um modo autônomo de apreender a tarefa comunicativa e responder as suas perguntas ( BRASIL. p. 94, 2018 ).
Como sujeito ativo, a \textbf{criança} \textbf{demonstra} \textbf{curiosidade} sobre o mundo físico e o mundo sociocultural.
Sente a necessidade de compreender seu próprio corpo, os animais, as plantas, as transformações da natureza, os fenômenos atmosféricos, os diferentes materiais e as diversas alternativas de sua manipulação.
No mundo sociocultural, as relações familiares e sociais permitem à \textbf{criança} conhecer e compreender a dinâmica da sociedade, os relacionamentos, a diversidade e as tradições.
Tanto as \textbf{crianças} das zonas urbanas, quanto as do campo e das comunidades tradicionais, \textbf{manipulam} objetos e materiais diversos para entender seu funcionamento e suas características.
Os conhecimentos lógico-matemáticos, como forma de comunicação, também se realizam a partir da \textbf{curiosidade} da \textbf{criança}, o que permite ao professor realizar experiências desafiadoras de analisar, observar, comparar, tomar decisões, resolver problemas, explorar ideias, construir e \textbf{testar} hipóteses e tirar conclusões.
Assim, é importante promover interações e \textbf{brincadeiras} que proporcionem ricas oportunidades de explorar os espaços, os objetos, formas, espessura, \textbf{texturas}, dimensões, cor, os tempos, as quantidades, as relações e as transformações para encontrar respostas às suas \textbf{curiosidades} e indagações.
Promover vivências práticas que cumpram os direitos de aprendizagem, fortaleçam a autonomia das \textbf{crianças}, contribuam no desenvolvimento de suas habilidades e na construção de conhecimentos sobre o mundo físico e sociocultural.
Espaços, tempos, quantidades, relações e transformações - CONVIVER com \textbf{crianças} e \textbf{adultos} e com eles criar estratégias para investigar o mundo social e natural, \textbf{demonstrando} atitudes positivas em relação a situações que envolvam diversidade étnico-racial, ambiental, de gênero, de língua, de religião. - \textbf{BRINCAR} com materiais e objetos cotidianos, associados a diferentes papéis ou cenas sociais, e com elementos da natureza que apresentam diversidade de formas, \textbf{texturas}, cheiros, cores, tamanhos, \textbf{pesos}, densidades, experimentando possibilidades de transformação. - PARTICIPAR de atividades que deem a oportunidade de observação de contextos diversos, \textbf{atentando} para características do ambiente e das histórias locais, utilizando ferramentas de conhecimento e instrumentos de \textbf{registro}, orientação e comunicação, como bússola, lanterna, lupa, máquina \textbf{fotográfica}, gravador, filmadora, projetor, computador e celular. - EXPLORAR e identificar as características do mundo natural e social, nomeando -as, reagrupando -as e ordenando -as, segundo critérios diversos. - EXPRESSAR suas observações, hipóteses e explicações sobre objetos, organismos vivos, fenômenos da natureza, características do ambiente, \textbf{personagens} e situações sociais, registrando -as por meio de desenhos, fotografias, gravações em áudio e vídeo, escritas e outras linguagens. - CONHECER -SE e construir sua identidade pessoal e cultural, identificando seus próprios interesses na relação com o mundo físico e social, apropriando -se dos costumes, das crenças e tradições de seus grupos de pertencimento e do patrimônio cultural, artístico, ambiental, científico e tecnológico.
Brasil, 2016 ( 2ª Versão da BNCC )

% matched lemmas: adulto, atentar, bebê, brincadeira, brincar, carro, conseguir, criança, curiosidade, demonstrar, desejo, fotográfico, infantil, manipular, parceiro, pequeno, personagem, peso, registro, testar, textura
\end{textsample}
